\begin{description}
\item[(a)] In a magnetic field, a charged particle moves along a circular
  path defined by the equality of centrifugal and magnetic forces:
  \[ mv^2/r = evB, \]
  \hfill(1 point)

  where $m$ is the mass, $v$ velocity, $r$ radius of the circle, $e$ charge
  of the particle, and $B$ magnetic flux density.

  For circular motion, $v = 2\pi r/T$, therefore $T = 2\pi m/eB$. The
  charged particle moving in harmonic motion (i.e.\ along the circular path
  with constant velocity) emits a wave of wavelength
  $\lambda = cT = 2\pi mc/eB$ and thus $B = 2\pi mc/e\lambda$.
  \hfill(1 point)

  Substituting the numerical values, including the mass and charge of the
  electron:
  \[ B = (2\pi\times9.109\times10^{-31}\times2.998\times10^{8})/
     (1.602\times10^{-19}\times6\times10^{-7})
     \approx 2\times10^{4}\ \mathrm{T}. \]
  \hfill(1 point)

  Since $\lambda$ is given to 1 s.f., the correct answer is 20\,kT, however
  accept 17.9\,kT or 18\,kT. More than 3 s.f.\ in the final answer loses
  points.
\item[(b)] In the plasma, besides electrons, only protons will be present in
  large quantities; protons will emit energy at a wavelength larger in
  proportion to the mass ratio, i.e.\ $1836\times$ larger (anything within
  1800--2000$\times$ $\implies$ $\lambda = 1.08$--$1.20$\,mm is acceptable).
  \hfill(2 points)
\end{description}
